% Generated by manuscript/build_manuscript.py.
\section{Discussion and Limitations}
\label{sec:discussion}


\textbf{Unified varied-horizon forecasting.} UVHF serves different request endpoints from a single prediction trajectory. Horizon-specific predictors optimize each endpoint independently and do not enforce CHPC across their overlapping outputs. HoriScope maps a shared history representation and future-step coordinates to one trajectory, whose prefixes answer requests of different lengths. The experiments in Sections~\ref{sec:horizon-specific-comparison} and~\ref{sec:one-model-comparison} show that this structural consistency coexists with strong forecasting accuracy, and Fig.~\ref{fig:accuracy-system-cost} demonstrates the deployment benefit of consolidating all horizon requests into one checkpoint.

\textbf{Output-side multi-scope generation.} Existing multi-scale encoders extract and combine historical dependencies at multiple input resolutions, where scale describes the granularity of history representation before decoding \citep{challu2023nhits,chen2024pathformer,wang2024timemixer}. HoriScope introduces a distinct output-side formulation in which candidate forecasts reuse history-conditioned latent states over different contiguous future extents. Each scope therefore specifies the sharing granularity of forecast generation, shifting multi-scale modeling from history representation to future construction. Evidence from the scope-wise ablations, the future-region analysis in Section 3 and the Fig.~\ref{fig:scope-allocation-behavior} diagnostics indicates that different future regions can benefit from different sharing extents. Target-Adaptive Allocation assigns step-specific preferences across these candidate forecast signals. Multi-scale encoders and HoriScope address separate yet complementary stages of the forecasting pipeline. Multi-scale encoders organize historical evidence, while HoriScope converts the resulting shared representation into forecasts at multiple future-sharing granularities.

\textbf{Methodological limitations.} The current formulation relies on a finite set of pre-specified contiguous scopes and future regions. This assumption may be restrictive when useful sharing patterns change at irregular locations, motivating data-adaptive region partitioning. Interactions among region forecasts are mediated through the shared History State and Future Coordinate. A dedicated cross-region forecasting mechanism could model dependencies that evolve across the future domain more explicitly. Multiple scope branches also increase decoder-side computation compared with a single-scope head, although shared encoding and unified serving reduce duplicated computation and storage. Future work should examine adaptive partitioning, cross-region interactions and more efficient multi-scope generation.
